\chapter{Clustering}
Clustering Algorithms

The clustering component of the RUFA system provides efficient spatial grouping of urban tree data to enable rapid analysis and visualization at various geographic scales. This chapter presents the evolution of clustering algorithms implemented, from basic K-means to advanced SuperClustering with Voronoi diagrams, demonstrating progressive performance improvements for handling California's 6.6 million tree inventory records.

\section{K-Means}

The initial clustering implementation utilized the traditional K-means algorithm for spatial grouping of tree coordinates. K-means partitions the dataset into k clusters by minimizing the within-cluster sum of squares (WCSS):

\begin{equation}
WCSS = \sum_{i=1}^{k} \sum_{x \in C_i} ||x - \mu_i||^2
\end{equation}

where $C_i$ represents cluster $i$, $\mu_i$ is the centroid of cluster $i$, and $x$ represents individual tree coordinates.

\textbf{Implementation Details:}
The basic K-means implementation processes tree coordinates using the Lloyd's algorithm with the following specifications:
\begin{itemize}
\item \textbf{Distance Metric:} Haversine distance for geographic coordinates
\item \textbf{Initialization:} K-means++ for optimal initial centroid selection
\item \textbf{Convergence Criteria:} Maximum 100 iterations or centroid movement < 0.001°
\item \textbf{Cluster Count:} Dynamic k selection based on zoom level (k = 50-500)
\end{itemize}

\textbf{Performance Characteristics:}
Initial benchmarking revealed significant computational overhead for large datasets:
\begin{itemize}
\item Average execution time: 2.3 seconds for 100,000 trees
\item Memory usage: 45MB for coordinate processing
\item Scalability limitation: O(nkt) complexity where n=points, k=clusters, t=iterations
\end{itemize}

The basic K-means approach provided adequate clustering quality but suffered from performance issues when processing the full California dataset, necessitating optimization strategies.

\section{K-Means with Caching}

To address performance limitations, the second iteration implemented comprehensive caching mechanisms for both intermediate computations and final cluster results.

\textbf{Caching Strategy:}
\begin{itemize}
\item \textbf{Centroid Cache:} Cached cluster centroids

\textbf{Performance Improvements:}
The caching implementation achieved significant performance gains:
\end{itemize}

\begin{itemize}
\item TBD
\item TBD
\end{itemize}

\section{Additional Approaches}

Recognizing the limitations of standard K-means for geographic data, several alternative clustering approaches were evaluated and implemented.

\subsection{Hierarchical Clustering}
Agglomerative hierarchical clustering was implemented to create multi-level tree groupings supporting zoom-dependent visualizations:

\begin{itemize}
\item \textbf{Linkage Criteria:} Ward linkage minimizing within-cluster variance
\item \textbf{Distance Matrix:} Sparse matrix optimization for geographic coordinates
\item \textbf{Dendrogram Pruning:} Dynamic tree cutting based on zoom level requirements
\end{itemize}

Performance analysis showed hierarchical clustering provided better multi-scale representations but with increased computational complexity O(n³), making it suitable only for smaller regional datasets.

\subsection{Grid-Based Clustering}
A custom grid-based approach was another approach for rapid initial clustering:

\begin{itemize}
\item \textbf{Grid Resolution:} Adaptive cell size based on tree density (0.001° - 0.01°)
\item \textbf{Spatial Index:} R-tree spatial indexing for efficient range queries
\item \textbf{Aggregation:} Statistical summaries computed per grid cell
\end{itemize}

Grid-based clustering achieved O(n) complexity but lacked the flexibility needed for variable-density urban environments.

\section{SuperClustering with Voronoi Diagrams}

This clustering implementation combines hierarchical clustering with Voronoi diagram generation to analyze urban tree data efficiently. The approach creates spatial boundaries at multiple zoom levels, enabling both detailed local analysis and broader regional insights.

\subsection{System Overview}

The system integrates two key technologies:
\begin{itemize}
\item \textbf{SuperCluster}: Hierarchical point clustering across zoom levels 0--10
\item \textbf{Voronoi Diagrams}: Spatial tessellation for regional boundary creation
\end{itemize}

This combination enables multi-scale spatial analysis with efficient point aggregation and real-time visualization capabilities.

\subsection{SuperCluster Configuration}

The clustering process uses pysupercluster with the following parameters:

\begin{lstlisting}[language=Python, caption=Basic SuperCluster Setup]
# Initialize clustering with key parameters
index = pysupercluster.SuperCluster(
    detector_points,
    min_zoom=0,
    max_zoom=10,
    radius=5,
    extent=512
)

# Process zoom levels
for zoom in range(5, 10):
    clusters = index.getClusters(bbox, zoom)
    if len(clusters) >= 4:
        voronoi_data = process_voronoi(clusters, zoom)
\end{lstlisting}

\textbf{Key Configuration Parameters:}
\begin{description}
\item[Zoom Range] 0 to 10 levels for hierarchical analysis
\item[Cluster Radius] 5-pixel aggregation radius
\item[Tile Extent] 512-pixel spatial indexing
\item[Coverage] Global bounding box (-180°, 90°) to (180°, -90°)
\end{description}

\subsection{Voronoi Diagram Generation}

For each zoom level with sufficient cluster points, Voronoi diagrams create spatial boundaries:

\begin{lstlisting}[language=Python, caption=Simplified Voronoi Processing]
def process_voronoi(cluster_points, zoom):
    # Generate Voronoi diagram
    vor = Voronoi(cluster_points)
    
    # Create polygons from Voronoi edges
    polygons = create_polygons_from_voronoi(vor)
    
    # Process each Voronoi cell
    for polygon in polygons:
        # Calculate area and tree metrics
        area_km = calculate_area(polygon)
        trees_in_cell = count_trees_in_polygon(polygon)
        
        # Compute diversity and density
        tree_density = trees_in_cell / area_km
        diversity_score = calculate_td_50(polygon)
        
        # Add to results with color mapping
        add_feature_to_geojson(polygon, tree_density, diversity_score)
\end{lstlisting}

\subsection{Spatial Analysis Metrics}

Each Voronoi cell generates comprehensive urban forestry metrics:

\begin{description}
\item[Tree Density] Trees per square kilometer within each cell
\item[Species Diversity] TD-50 index measuring ecological diversity
\item[Color Mapping] Normalized visualization based on density values
\item[Geometric Validation] Ensures valid polygon output for mapping
\end{description}

\subsection{Point Assignment Algorithm}

Trees are efficiently assigned to Voronoi cells using spatial containment:

\begin{lstlisting}[language=Python, caption=Point-in-Polygon Assignment]
def assign_trees_to_cells(points, voronoi_polygons):
    for polygon in voronoi_polygons:
        trees_in_cell = []
        for tree_point in points:
            if polygon.contains(tree_point):
                trees_in_cell.append(tree_point)
        
        # Calculate metrics for this cell
        process_cell_metrics(polygon, trees_in_cell)
\end{lstlisting}

\subsection{TD-50 Diversity Calculation}

The TD-50 index measures species diversity within each Voronoi cell:

\begin{lstlisting}[language=Python, caption=Regional Diversity Assessment]
def calculate_td_50(species_counts):
    if not species_counts:
        return 0
    
    # Sort species by abundance
    sorted_counts = sorted(species_counts, reverse=True)
    total_trees = sum(sorted_counts)
    
    # Find species needed for 50% of trees
    cumulative = 0
    for i, count in enumerate(sorted_counts):
        cumulative += count
        if cumulative >= total_trees * 0.5:
            return i + 1
\end{lstlisting}

\subsection{Performance Characteristics}

The SuperClustering implementation demonstrates significant computational efficiency:

\textbf{Algorithmic Complexity:}
\begin{itemize}
\item SuperCluster initialization: $O(n \log n)$
\item Voronoi generation: $O(n \log n)$ using Fortune's algorithm
\item Point assignment: $O(kn)$ where $k$ is the number of cells
\item Overall complexity: $O(n \log n + kn)$
\end{itemize}

\textbf{Multi-Scale Benefits:}
\begin{itemize}
\item Higher zoom levels provide detailed local analysis
\item Lower zoom levels enable regional pattern recognition
\item Cached results allow rapid zoom-level switching
\item Hierarchical structure supports interactive exploration
\end{itemize}

\textbf{Output Features:}
\begin{itemize}
\item Valid GeoJSON format for web mapping integration
\item Color-coded visualization based on tree density
\item Species diversity metrics for ecological assessment
\item Scalable performance across California's metropolitan regions
\end{itemize}

This approach successfully combines the hierarchical efficiency of SuperCluster with the spatial analysis power of Voronoi diagrams, enabling comprehensive urban forestry assessment at multiple scales.